% Os 496 pregões do traço (jul/2019 a jun/2021) passando pelo modelo: quantos o
% ARIMA mandou para cada lado, quantos o GARCH vetou e o que o robô fez.
% Contagens: dados/resultados/robo_em_acao_2026-10-30/valores.tex (\pR...),
% recontadas a partir do dia_a_dia.csv em 06/11/2026 (240 + 56 + 64 + 136).
\documentclass[border=6pt]{standalone}
% Preâmbulo comum das figuras de ia/relatorios/modelo-completo e garch-funcionando
% (06/11/2026). Nenhum número é digitado aqui: vêm de valores.tex (\vX para o
% texto, \pX para coordenada) e dos CSVs de
% dados/resultados/modelo_completo_2026-11-06/ (ferramentas/ModeloCompleto.mq5)
% e dados/resultados/robo_em_acao_2026-10-30/ (traço do robô).
\usepackage[T1]{fontenc}
\usepackage{amsmath}
\usepackage{mathptmx}
\usepackage{pgfplots}
\usepgfplotslibrary{dateplot,fillbetween}
\usetikzlibrary{arrows.meta,positioning,calc,decorations.pathreplacing}
\pgfplotsset{compat=1.18}
\definecolor{cmod}{RGB}{0,90,181}     % ARIMA, ajuste
\definecolor{cdes}{RGB}{220,50,32}    % GARCH acima do limiar, veto
\definecolor{cdad}{RGB}{120,120,120}  % dado observado
\definecolor{cok}{RGB}{0,140,70}      % compra liberada
\definecolor{cgar}{RGB}{230,140,0}    % GARCH, limiar
\newcommand{\arqm}[1]{../../../dados/resultados/modelo_completo_2026-11-06/#1}
\newcommand{\arqr}[1]{../../../dados/resultados/robo_em_acao_2026-10-30/#1}
\input{\arqm{valores.tex}}
\input{\arqr{valores.tex}}
\newcommand{\passo}[1]{\tikz[baseline=(n.base)]\node[circle,fill=black!75,text=white,inner sep=1pt,font=\footnotesize\bfseries](n){#1};}
\pgfplotsset{
  mini/.style={table/col sep=comma, scale only axis, unbounded coords=jump,
    tick label style={font=\scriptsize}, label style={font=\scriptsize},
    grid=major, grid style={gray!20}, every axis plot/.append style={line width=0.7pt},
    /pgf/number format/use comma, /pgf/number format/1000 sep={.}},
  datas/.style={date coordinates in=x},
}
\tikzset{
  caixa/.style={draw=black!45, rounded corners=3pt, line width=0.6pt},
  titulo/.style={anchor=north west, font=\small\bfseries, inner sep=0pt},
  texto/.style={anchor=north west, font=\footnotesize, align=left, inner sep=0pt},
  seta/.style={-{Stealth[length=2.6mm]}, line width=1pt, black!55},
}

\begin{document}
\begin{tikzpicture}[x=1cm,y=1cm]
\pgfmathsetmacro{\k}{16/\pRN}
\pgfmathsetmacro{\xa}{\k*\pRCompras}                 % fim de "alta liberada"
\pgfmathsetmacro{\xb}{\xa+\k*\pRVetoCompra}          % fim de "alta vetada" = fim da alta
\pgfmathsetmacro{\xc}{\xb+\k*\pRVetoVenda}           % fim de "queda vetada"
% faixa: \faixa{x0}{x1}{y}{cor}{texto}
\newcommand{\seg}[5]{\fill[#4] (#1,#3) rectangle (#2,#3-0.8); \draw[white, line width=1pt] (#1,#3) rectangle (#2,#3-0.8);
  \node[font=\footnotesize, align=center] at ({(#1+#2)/2},#3-0.4) {#5};}
\newcommand{\fita}[4]{\fill[#4, opacity=0.35] (#1,#3) rectangle (#2,#3-0.55);}
\newcommand{\rotulo}[3]{\node[anchor=east, font=\footnotesize\bfseries, align=right, text=#3] at (-0.25,#1-0.4) {#2};}

\node[anchor=south west, font=\large\bfseries] at (-3.6,0.55) {O caminho de \pRN{} pregões pelo modelo (\vRInicio{} a \vRFim)};

% linha 0
\rotulo{0}{pregões}{black}
\seg{0}{16}{0}{black!12}{\pRN{} pregões, um sinal novo a cada abertura}
\fita{0}{\xb}{-0.8}{cmod!50}\fita{\xb}{16}{-0.8}{black!30}

% linha 1: ARIMA
\rotulo{-1.35}{ARIMA\\(direção)}{cmod}
\seg{0}{\xb}{-1.35}{cmod!35}{previsão de alta: \pRComprasA}
\seg{\xb}{16}{-1.35}{black!18}{previsão de queda: \pRVendasA}
\fita{0}{\xa}{-2.15}{cok!60}\fita{\xa}{\xc}{-2.15}{cdes!60}\fita{\xc}{16}{-2.15}{black!30}

% linha 2: GARCH
\rotulo{-2.7}{GARCH\\(risco)}{cgar!85!black}
\seg{0}{\xa}{-2.7}{cok!30}{liberada: \pRCompras}
\seg{\xa}{\xb}{-2.7}{cdes!35}{vetada:\\\pRVetoCompra}
\seg{\xb}{\xc}{-2.7}{cdes!35}{vetada:\\\pRVetoVenda}
\seg{\xc}{16}{-2.7}{black!18}{sem veto: \pRVendas}
\fita{0}{\xa}{-3.5}{cok!60}\fita{\xa}{16}{-3.5}{black!30}

% linha 3: ação
\rotulo{-4.05}{o robô}{black}
\pgfmathtruncatemacro{\fora}{\pRN-\pRCompras}
\seg{0}{\xa}{-4.05}{cok!45}{comprado: \pRCompras}
\seg{\xa}{16}{-4.05}{black!25}{fora do mercado (zera a posição): \fora}

% chaves de leitura
\draw[decorate, decoration={brace, amplitude=5pt, mirror}, black!60] (\xa,-4.95) -- node[below=6pt, font=\footnotesize, align=center, text width=6.6cm]
  {os \pRVetos{} vetos (\pRVetosPct\% dos pregões) tiram o robô do mercado nos dois sentidos} (\xc,-4.95);
\node[anchor=north west, font=\footnotesize, align=left, text width=19.2cm] at (-3.6,-6.25) {Como ler: cada pregão
  desce pelas linhas. O ARIMA decide o lado, o GARCH decide se o robô pode ficar exposto. O GARCH só retira: nenhum pregão
  com previsão de queda vira compra, e a versão sem GARCH, com os mesmos sinais, ficaria comprada nos \pRComprasA{} de alta.
  Como o robô só compra, previsão de queda e veto levam ao mesmo lugar, fora do mercado. As larguras são proporcionais
  ao número de pregões. Configuração ARIMA-GARCH Sharpe (filtro, percentil 70, zera no veto), PETR4.};
\end{tikzpicture}
\end{document}
